% =============================================================================
\section{Related Work}
\label{sec:related-work}
% =============================================================================

This section surveys four streams of literature that inform the PURECORTEX
tokenomic design: bonding curves and automated market makers, vote-escrowed
tokenomics, agent tokenization platforms, and game-theoretic mechanism design
for decentralized systems.

\subsection{Bonding Curves and Continuous Token Models}
\label{subsec:rw-bonding}

The concept of automated price discovery through mathematical functions
originates with Hanson's logarithmic market scoring rule for prediction
markets~\citep{hanson2003combinatorial}. Hanson demonstrated that a cost
function derived from a convex potential can simultaneously serve as a market
maker and an information aggregation device. The Bancor protocol generalized
this insight to arbitrary token issuance, introducing continuous liquidity
through a reserve-backed bonding curve parameterized by a connector
weight~\citep{bancor2017}. Subsequent work by Angeris et
al.~\citep{angeris2020improved} provided rigorous analysis of constant-product
market makers (the Uniswap model), establishing conditions under which
arbitrageurs drive prices toward external reference values.

Our quadratic bonding curve (Section~\ref{sec:bonding-curve}) builds on this
lineage but introduces two innovations absent from prior work: a deterministic
graduation mechanism that atomically transitions liquidity from the bonding
curve to a decentralized exchange, and a mandatory creator vesting schedule
that conditions token release on elapsed time rather than price milestones.

\subsection{Vote-Escrowed Tokenomics}
\label{subsec:rw-ve}

The vote-escrowed (ve) model was pioneered by Curve Finance's
veCRV~\citep{curvefinance2020, egorov2019stableswap}, wherein users lock CRV
tokens for up to four years to receive governance weight proportional to the
remaining lock duration. This mechanism simultaneously reduces circulating
supply, aligns governance power with long-term commitment, and distributes
protocol revenue to lockers. Balancer subsequently adapted the model as
veBAL~\citep{balancer2022vebal}, adding a gauge-based emission allocation
system.

Two limitations of existing ve implementations motivate our design. First,
veCRV was introduced approximately two years after CRV launch, disadvantaging
early participants who could not lock tokens from inception. Second, the linear
decay of governance weight creates a ``cliff problem'' wherein users must
repeatedly re-lock to maintain influence, introducing gas costs and
attentional overhead. The \veCORTEX{} mechanism
(Section~\ref{sec:vecortex-staking}) addresses the first limitation by
launching contemporaneously with the token, and mitigates the second through a
boost multiplier that rewards sustained participation.

\subsection{Agent Tokenization Platforms}
\label{subsec:rw-agent}

The tokenization of AI agents---the process of creating tradeable digital
assets representing ownership stakes in or usage rights to autonomous
agents---has emerged as a distinct category within decentralized infrastructure.
The most prominent exemplar, which we refer to as \emph{Platform~A} for
neutrality, implements the following design:

\begin{itemize}
  \item A fixed supply of $10^9$ tokens per agent, with 1\% allocated to the
    creator at launch and no vesting schedule.
  \item A flat 1\% trading fee on all buy and sell transactions, split between
    the creator and the protocol.
  \item Initial liquidity bootstrapped through a bonding curve, with LP tokens
    locked upon migration to a decentralized exchange.
  \item No on-chain governance mechanism; protocol parameters are controlled by
    the development team.
  \item No composability framework; agents operate as isolated services.
\end{itemize}

While Platform~A demonstrated product-market fit and catalyzed the agent token
category, its design exhibits the five structural deficiencies enumerated in
Section~\ref{subsec:problem}. The absence of creator vesting has been
empirically associated with rapid creator exits; the flat fee structure fails
to adapt to volatility regimes; and the lack of governance creates platform
risk.

\subsection{Game Theory and Mechanism Design in DeFi}
\label{subsec:rw-game-theory}

Roughgarden~\citep{roughgarden2020transaction} established a formal framework
for analyzing transaction fee mechanisms as games between users and block
producers, proving that EIP-1559-style mechanisms are incentive-compatible
under certain conditions. Daian et al.~\citep{daian2020flash} identified
miner extractable value (MEV) as a fundamental game-theoretic challenge in
decentralized exchanges, demonstrating that frontrunning and sandwich attacks
constitute rational strategies under existing fee structures.

In the broader mechanism design literature,
Myerson~\citep{myerson1981optimal} established revenue-optimal auction theory,
while Spence~\citep{spence1973job} introduced the signaling model that we
adapt in Section~\ref{sec:game-theory} to analyze bonding curve participation
as a quality signal. Zargham et al.~\citep{zargham2020} proposed cadCAD-based
simulation frameworks for economic games in token ecosystems, providing
methodological grounding for our simulation approach.

The intersection of game theory and governance has been explored in the context
of liquid democracy and delegative voting~\citep{fudenberg1991game}. Our
governance model (Section~\ref{sec:governance}) extends this literature by
introducing an AI Senator agent that participates in the proposal process but
is constitutionally prohibited from voting, creating a novel separation of
powers within a token-governed system.

\subsection{Algorand as Settlement Layer}
\label{subsec:rw-algorand}

The Algorand blockchain~\citep{algorand2019, chen2019algorand} provides
several properties that are structurally advantageous for agent tokenomics:
immediate transaction finality (no probabilistic confirmation), native asset
support through the Algorand Standard Asset (ASA) model, atomic transaction
groups that enable multi-step operations without reentrancy risk, and
LogicSig-based programmable accounts that can enforce time-locked LP
positions without trusted intermediaries.

These properties eliminate entire classes of vulnerabilities present in
EVM-based implementations (see Section~\ref{sec:security}) and enable the
atomic graduation mechanism described in Section~\ref{sec:bonding-curve}.
